\hnsection{LINEAR REPRESENTATIVES AND GLOBAL LIE GROUPS}

None of the works which we have just mentioned made any honest attempt to define and study global Lie groups. It was H. Weyl who made the first move in this direction. He was inspired by two theories, which until then had developed independently: that of linear representations of complex semi-simple Lie algebras, due to E. Cartan, and that of linear representations of finite groups, due to Frobenius and which had just been transposed to the orthogonal group by I. Schur, using an idea of Hurwitz. The later had shown {[}17{]} how to form invariants for the orthogonal group or unitary group by replacing the operation of the mean on a finite group by an integration with respect to an invariant measure. He had also noted that, by applying this method to the unitary group, invariants are obtained for the general linear group, the first example of the ``unitarian trick''. In 1924, I. Schur {[}20{]} used this procedure to show the complete reducibility of the representations of the orthogonal group \(\mathbf{O}(n)\) and the unitary group \(\mathbf{U}(n)\) by constructing an invariant non-degenerate positive definite Hermitian form; he deduced, by the ``unitarian trick'', the complete reducibility of the holomorphic representations of \(\mathbf{O}(n,\mathbf{C})\) and \(\mathbf{SL}(n,\mathbf{C})\) , established orthogonality relations for the characters of \(\mathbf{O}(n)\) and \(\mathbf{U}(n)\) and determined the characters of \(\mathbf{O}(n)\) . H. Weyl immediately extended this method to complex semi-simple Lie algebras {[}21{]}. Given such an algebra g, he showed that it has a ``compact real form'' (which amounts to saying that is obtained by extension of scalars from R to C from an algebra \(g_{0}\) over R whose adjoint group \(G_{0}\) is compact). Further, he showed that the fundamental group of \(G_{0}\) is finite and hence that the universal covering‡ of \(G_{0}\) is compact. He deduced, by a suitable adaptation of Schur's procedure, the complete reducibility of the representations of g and also gave, by a global method, the determination of the characters of the representations of g. In a letter to I. Schur (Sitzungsber, Berlin, 1924, pp.~338--343), H. Weyl summarized the results of Cartan not known to Schur (cf.~{[}20{]}, p.~299, note at the bottom of the page) and compared the two approaches: Cartan's method gave all the holomorphic representations of the simply connected group with Lie algebra g; in the case of the orthogonal group, representations were thus obtained of a double covering (called later the spinor group), which had escaped Schur; on the other hand, Schur's method had the advantage of showing the complete reducibility and giving the characters explicitly.

After the work of H. Weyl, E. Cartan adopted an openly global approach in his research on symmetric spaces and Lie groups. This approach was the basis of his exposition of 1930 ({[}12{]}, vol.~I₂, pp.~1165--1225) of the theory of ``finite continuous'' groups. In particular we find here the first proof of the global version of the 3rd fundamental theorem (the existence of a Lie group with given Lie algebra); Cartan also showed that every closed subgroup of a real Lie group is a Lie group (Chapter III, § 8, no. 2, Theorem 2) which generalized a result of J. von Neumann on closed subgroups of the linear group {[}23{]}. In this Memoir, von Neumann showed also that every continuous representation of a semi-simple group is real analytic.

After these works, the theory of Lie groups in the ``classical'' sense (that is in finite dimension over R or C) was more or less complete in its essentials. The first detailed exposition of it was given by Pontrjagin in his book on topological groups {[}36{]}; there he followed an approach quite close to that of Lie, but carefully distinguished between the local and the global. This was followed by Chevalley's book {[}38{]} which also contains the first systematic discussion of the theory of analytic manifolds and exterior differential calculus; the ``infinitesimal transformations'' of Lie appear there as vector fields and the Lie algebra of a Lie group G is identified with the space of left invariant fields on G. He leaves out any discussion on ``group germs'' and ``transformation groups''.

